\chapter{1985年全国卷（理）}

\section{单选题}

\begin{problem}

如果正方体 \(ABCD - A'B'C'D'\) 的棱长为 \(a\)，那么四面体 \(A' - ABD\) 的体积是（\quad）

\choices
    {\(\frac{a^{3}}{2}\)}
    {\(\frac{a^3}{3}\)}
    {\(\frac{a^3}{4}\)}
    {\(\frac{a^3}{6}\)}

\begin{answer}
D
\end{answer}

\begin{solution}
三角形 \(ABD\) 是正方形 \(ABCD\) 的一半，所以其面积为 \(\frac{a^2}{2}\). 点 \(A'\) 到平面 \(ABD\) 的距离就是正方体的棱长 \(a\)，因此四面体的体积为 \(\frac{1}{3}\cdot\frac{a^2}{2}\cdot a=\frac{a^3}{6}\)，故选 D.
\end{solution}
\end{problem}

\begin{problem}

\(\tan x = 1\) 是 \(x = \frac{5}{4}\pi\) 的（\quad）

\choices
    {必要条件}
    {充分条件}
    {充分必要条件}
    {既不充分又不必要的条件}

\begin{answer}
A
\end{answer}

\begin{solution}
当 \(x=\frac{5\pi}{4}\) 时，\(\tan x=1\)，所以 \(x=\frac{5\pi}{4}\) 能推出 \(\tan x=1\)，即 \(\tan x=1\) 是该条件的必要条件. 反过来，\(\tan x=1\) 时还可能有 \(x=\frac{\pi}{4}+k\pi\)，并不一定等于 \(\frac{5\pi}{4}\)，所以它不是充分条件，故选 A.
\end{solution}
\end{problem}

\begin{problem}

在下面给出的函数中，哪一个函数既是区间 \(\left(0, \frac{\pi}{2}\right)\) 上的增函数又是以 \(\pi\) 为周期的偶函数（\quad）

\choices
    {\(y = x^{2}\)（\(x\in \R\)）}
    {\(y = |\sin x|\)（\(x\in \R\)）}
    {\(y = \cos 2x\)（\(x\in \R\)）}
    {\(y = \e^{\sin 2x}\)（\(x \in \R\)）}

\begin{answer}
B
\end{answer}

\begin{solution}
函数 \(y=x^2\) 虽为偶函数，但 \((x+\pi)^2\) 不恒等于 \(x^2\)，所以不是以 \(\pi\) 为周期的函数. 函数 \(y=|\sin x|\) 满足 \(|\sin(-x)|=|\sin x|\)，是偶函数；又有 \(|\sin(x+\pi)|=|\sin x|\)，以 \(\pi\) 为周期. 在 \(\left(0,\frac{\pi}{2}\right)\) 内，\(\sin x\) 为增函数且为正数，因此 \(|\sin x|=\sin x\) 为增函数. 函数 \(y=\cos 2x\) 在该区间内递减，函数 \(y=\e^{\sin 2x}\) 也不是偶函数，因为 \(\e^{\sin(-2x)}=\e^{-\sin 2x}\) 一般不等于 \(\e^{\sin 2x}\). 故选 B.
\end{solution}
\end{problem}

\begin{problem}

极坐标方程 \(\rho = a\sin \theta\)（\(a > 0\)）的图象是（\quad）

\choices
    {\choicebitmap[width=0.15\paperwidth]{img/1985/national_science/q4_choiceA.png}}
    {\choicebitmap[width=0.15\paperwidth]{img/1985/national_science/q4_choiceB.png}}
    {\choicebitmap[width=0.15\paperwidth]{img/1985/national_science/q4_choiceC.png}}
    {\choicebitmap[width=0.15\paperwidth]{img/1985/national_science/q4_choiceD.png}}

\begin{answer}
C
\end{answer}

\begin{solution}
由极坐标与直角坐标的关系 \(x=\rho\cos\theta\)，\(y=\rho\sin\theta\) 得 \(\rho^2=x^2+y^2\)，且 \(\rho\sin\theta=y\). 将 \(\rho=a\sin\theta\) 两边乘以 \(\rho\)，得到 \(x^2+y^2=ay\)，即 \(x^2+\left(y-\frac{a}{2}\right)^2=\left(\frac{a}{2}\right)^2\). 因此图象是圆心为 \(\left(0,\frac{a}{2}\right)\)、半径为 \(\frac{a}{2}\) 且与 \(x\) 轴在原点相切的圆，故选 C.
\end{solution}
\end{problem}

\begin{problem}

用 \(1\)，\(2\)，\(3\)，\(4\)，\(5\) 这五个数字，可以组成比 \(20000\) 大，并且百位数不是数字 \(3\) 的没有重复数字的五位数，共有（\quad）

\choices
    {\(96\text{ 个}\)}
    {\(78\text{ 个}\)}
    {\(72\text{ 个}\)}
    {\(64\text{ 个}\)}

\begin{answer}
B
\end{answer}

\begin{solution}
五位数大于 \(20000\)，首位只能是 \(2\)，\(3\)，\(4\) 或 \(5\). 首位为 \(2\) 时，后四位共有 \(4!=24\) 种排列，其中百位为 \(3\) 的有 \(3!=6\) 种，所以有 \(24-6=18\) 种. 首位为 \(3\) 时，百位不可能再为 \(3\)，有 \(4!=24\) 种. 首位为 \(4\) 或 \(5\) 时，各有 \(4!-3!=18\) 种. 总数为 \(18+24+18+18=78\)，故选 B.
\end{solution}
\end{problem}

\section{解答题}

\begin{problem}

求方程 \(2\sin \left(x + \frac{\pi}{6}\right) = 1\) 解集.

\begin{solution}
原方程等价于 \(\sin\left(x+\frac{\pi}{6}\right)=\frac{1}{2}\). 因此 \(x+\frac{\pi}{6}=\frac{\pi}{6}+2k\pi\) 或 \(x+\frac{\pi}{6}=\frac{5\pi}{6}+2k\pi\)，其中 \(k\in\Z\). 分别解得 \(x=2k\pi\) 或 \(x=\frac{2\pi}{3}+2k\pi\)，所以解集为 \(\left\{2k\pi\mid k\in\Z\right\}\cup\left\{\frac{2\pi}{3}+2k\pi\mid k\in\Z\right\}\).
\end{solution}
\end{problem}

\begin{problem}

设 \(|a|\leqslant 1\)，求 \(\arccos a+\arccos\left(-a\right)\) 的值.

\begin{solution}
令 \(\alpha=\arccos a\)，则 \(\alpha\in[0,\pi]\)，且 \(a=\cos\alpha\). 于是 \(-a=-\cos\alpha=\cos\left(\pi-\alpha\right)\)，并且 \(\pi-\alpha\in[0,\pi]\). 由于反余弦函数的值域为 \([0,\pi]\)，故 \(\arccos(-a)=\pi-\alpha\). 因此 \(\arccos a+\arccos(-a)=\alpha+\pi-\alpha=\pi\).
\end{solution}
\end{problem}

\begin{problem}

求曲线 \(y^{2}=-16x+64\) 的焦点.

\begin{solution}
将方程化为 \(y^2=-16(x-4)\)，与抛物线标准方程 \(y^2=4p(x-h)\) 比较，得 \(h=4\)，\(p=-4\). 所以顶点为 \((4,0)\)，焦点为 \((h+p,0)=(0,0)\).
\end{solution}
\end{problem}

\begin{problem}

设 \((3x-1)^{6}=a_{6}x^{6}+a_{5}x^{5}+a_{4}x^{4}+a_{3}x^{3}+a_{2}x^{2}+a_{1}x+a_{0}\)，求 \(a_{6}+a_{5}+a_{4}+a_{3}+a_{2}+a_{1}+a_{0}\) 的值.

\begin{solution}
令 \(x=1\) 代入已知恒等式，左边为 \((3-1)^6=2^6=64\)，右边为 \(a_6+a_5+a_4+a_3+a_2+a_1+a_0\). 因此所求值为 \(64\).
\end{solution}
\end{problem}

\begin{problem}

设函数 \(f(x)\) 的定义域是 \([0,1]\)，求函数 \(f(x^2)\) 的定义域.

\begin{solution}
要使 \(f(x^2)\) 有意义，必须有 \(x^2\in[0,1]\)，即 \(0\leqslant x^2\leqslant1\). 这等价于 \(-1\leqslant x\leqslant1\)，所以函数 \(f(x^2)\) 的定义域为 \([-1,1]\).
\end{solution}
\end{problem}

\begin{problem}

解方程：\(\log_4(3-x)+\log_{0.25}(3+x)=\log_4(1-x)+\log_{0.25}(2x+1)\).

\begin{solution}
首先由各对数的真数为正数，得到 \(3-x>0\)，\(3+x>0\)，\(1-x>0\)，\(2x+1>0\)，所以定义域为 \(-\frac{1}{2}<x<1\). 又因为 \(\log_{0.25}u=-\log_4u\)，原方程化为 \(\log_4(3-x)-\log_4(3+x)=\log_4(1-x)-\log_4(2x+1)\). 利用对数的减法公式并移项，得 \(\log_4\left((3-x)(2x+1)\right)=\log_4\left((3+x)(1-x)\right)\). 由于底数 \(4>1\)，对数函数是单调函数，故 \((3-x)(2x+1)=(3+x)(1-x)\). 展开得 \(-2x^2+5x+3=3-2x-x^2\)，即 \(x(7-x)=0\)，所以候选值为 \(x=0\) 或 \(x=7\). 其中 \(x=7\) 不在定义域内，只有 \(x=0\) 满足原方程，故解为 \(x=0\).
\end{solution}
\end{problem}

\begin{problem}

解不等式：\(\sqrt{2x+5}>x+1\).

\begin{solution}
由根式有意义，先得 \(x\geqslant-\frac{5}{2}\). 当 \(-\frac{5}{2}\leqslant x<-1\) 时，\(x+1<0\)，而 \(\sqrt{2x+5}\geqslant0\)，所以不等式恒成立. 当 \(x\geqslant-1\) 时，两边均非负，可以平方，得到 \(2x+5>(x+1)^2\)，即 \(x^2<4\)，所以 \(-2<x<2\). 与 \(x\geqslant-1\) 取交集得 \(-1\leqslant x<2\). 合并两种情况，原不等式的解集为 \(\left[-\frac{5}{2},2\right)\).
\end{solution}
\end{problem}

\begin{problem}

如图，设平面 \(AC\) 和 \(BD\) 相交于 \(BC\)，它们所成的一个二面角为 \(45^{\circ}\)，\(P\) 为平面 \(AC\) 内的一点，\(Q\) 为平面 \(BD\) 内的一点. 已知直线 \(MQ\) 是直线 \(PQ\) 在平面 \(BD\) 内的射影，并且 \(M\) 在 \(BC\) 上. 又设 \(PQ\) 与平面 \(BD\) 所成的角为 \(\beta\)，\(\angle CMQ=\theta\)（\(0^{\circ}<\theta<90^{\circ}\)），线段 \(PM\) 的长为 \(a\)，求线段 \(PQ\) 的长.

\bitmapfigure[width=0.52\linewidth,trim=4.35cm 6.6cm 4.15cm 1.7cm,clip]{img/1985/national_science/q13_fig1.png}

\begin{solution}
从 \(P\) 向平面 \(BD\) 作垂线，垂足为 \(R\). 因为 \(MQ\) 是 \(PQ\) 在平面 \(BD\) 内的射影，所以 \(R\) 在直线 \(MQ\) 上. 过 \(R\) 向交线 \(BC\) 作垂线，垂足为 \(N\). 由三垂线定理，\(PN\perp BC\)，而 \(RN\perp BC\)，所以 \(\angle PNR\) 是两平面所成二面角的平面角，从而 \(\angle PNR=45^{\circ}\). 又因 \(PR\perp BD\)，且 \(RN\subset BD\)，所以 \(PR\perp RN\)，三角形 \(PNR\) 在 \(R\) 处直角，故 \(PR=RN\).

令 \(\varphi=\angle PMR\). 在直角三角形 \(PMR\) 中，\(PM=a\)，所以 \(MR=a\cos\varphi\)，\(PR=a\sin\varphi\). 由于 \(R\) 在 \(MQ\) 上，\(\angle CMR=\angle CMQ=\theta\). 在直角三角形 \(MNR\) 中，\(NR=MR\sin\theta=a\cos\varphi\sin\theta\). 结合 \(PR=NR\)，得到 \(a\sin\varphi=a\cos\varphi\sin\theta\)，即 \(\tan\varphi=\sin\theta\).

在三角形 \(PMQ\) 中，\(\angle PMQ=\varphi\)，而 \(\angle PQM=\beta\)，后者是直线 \(PQ\) 与平面 \(BD\) 所成的角. 由正弦定理，\(\frac{PQ}{\sin\varphi}=\frac{PM}{\sin\beta}\)，所以 \(PQ=\frac{a\sin\varphi}{\sin\beta}\). 由 \(\tan\varphi=\sin\theta\) 且 \(0<\varphi<90^{\circ}\)，有 \(\sin\varphi=\frac{\sin\theta}{\sqrt{1+\sin^2\theta}}\). 因此 \(PQ=\frac{a\sin\theta}{\sin\beta\sqrt{1+\sin^2\theta}}\).
\end{solution}
\end{problem}

\begin{problem}

设 \(O\) 为复平面的原点，\(Z_{1}\) 和 \(Z_{2}\) 为复平面内的两动点，并且满足：

\begin{enumerate}
\item \(Z_{1}\) 和 \(Z_{2}\) 所对应的复数的辐角分别为定值 \(\theta\) 和 \(-\theta\)（\(0<\theta<\frac{\pi}{2}\)）；
\item \(\triangle OZ_{1}Z_{2}\) 的面积为定值 \(S\)，求 \(\triangle OZ_{1}Z_{2}\) 的重心 \(Z\) 所对应的复数的模的最小值.
\end{enumerate}

\begin{solution}
设 \(Z_1\)，\(Z_2\) 和重心 \(Z\) 所对应的复数分别为 \(z_1\)，\(z_2\) 和 \(z\)，并设 \(|z_1|=r_1>0\)，\(|z_2|=r_2>0\). 由辐角条件，\(z_1=r_1(\cos\theta+\i\sin\theta)\)，\(z_2=r_2(\cos\theta-\i\sin\theta)\). 三角形面积条件给出 \(S=\frac{1}{2}r_1r_2\sin 2\theta\)，所以 \(r_1r_2=\frac{S}{\sin\theta\cos\theta}\).

由于重心对应的复数满足 \(3z=z_1+z_2\)，有 \(3z=(r_1+r_2)\cos\theta+(r_1-r_2)\i\sin\theta\). 因此
\[
|3z|^2=(r_1+r_2)^2\cos^2\theta+(r_1-r_2)^2\sin^2\theta=(r_1-r_2)^2+4r_1r_2\cos^2\theta
\]
由 \((r_1-r_2)^2\geqslant0\)，得 \(|3z|^2\geqslant4r_1r_2\cos^2\theta=4S\cot\theta\). 当且仅当 \(r_1=r_2\) 时取等号，且此时满足给定的面积条件，因此 \(|z|\) 的最小值为 \(\frac{2}{3}\sqrt{S\cot\theta}\).
\end{solution}
\end{problem}

\begin{problem}

已知两点 \(P\left(-2,2\right)\)，\(Q\left(0,2\right)\) 以及一条直线 \(L:y=x\)，设长为 \(\sqrt{2}\) 的线段 \(AB\) 在直线 \(L\) 上移动，如图，求直线 \(PA\) 和 \(QB\) 的交点 \(M\) 的轨迹方程. （要求把结果写成普通方程）

\bitmapfigure[width=0.65\linewidth,trim=3.5cm 4.2cm 1.5cm 1.4cm,clip]{img/1985/national_science/q15_fig1.png}

\begin{solution}
设 \(A=(t,t)\)，\(B=(t+1,t+1)\)，其中 \(t\in\R\)，这样 \(AB=\sqrt{2}\). 设交点为 \(M=(x,y)\). 由直线 \(PA\) 经过 \((-2,2)\) 和 \((t,t)\)，直线 \(QB\) 经过 \((0,2)\) 和 \((t+1,t+1)\)，分别得到
\[
t(y-x-4)+2(x+y)=0
\]
\[
t(y-x-2)+x+y-2=0
\]
当 \(t=0\) 时，上述两式分别化为 \(x+y=0\) 和 \(x+y=2\)，故两直线平行且无交点. 以下设 \(t\neq0\)，令 \(u=x-y\)，\(v=x+y\)，上面两式化为 \(2v=t(u+4)\) 和 \(v=t(u+2)+2\). 消去 \(v\) 得 \(tu=-4\)，再代回得 \(v=2t-2\). 于是 \(u(v+2)+8=0\). 代回 \(u=x-y\)，\(v=x+y\)，可得
\[
x^2-y^2+2x-2y+8=0
\]
反过来，满足该方程的点必有 \(x-y\neq0\)，取 \(t=-\frac{4}{x-y}\) 即可由上述两式得到对应的交点，因此这就是交点 \(M\) 的轨迹普通方程.
\end{solution}
\end{problem}

\begin{problem}

设 \(a_{n}=\sqrt{1\cdot 2}+\sqrt{2\cdot 3}+\cdots+\sqrt{n(n+1)}\)（\(n=1\)，\(2\)，\(\cdots\)）.

\begin{enumerate}
\item 证明：不等式 \(\frac{n(n+1)}{2}<a_n<\frac{(n+1)^2}{2}\) 对所有的正整数 \(n\) 都成立；
\item 设 \(b_n=\frac{a_n}{n(n+1)}\)（\(n=1\)，\(2\)，\(\cdots\)），用定义证明：\(\displaystyle\lim_{n\to\infty}b_n=\frac{1}{2}\).
\end{enumerate}

\begin{solution}
\begin{enumerate}
\item 对任意正整数 \(k\)，有 \(k^2<k(k+1)<\left(k+\frac{1}{2}\right)^2\)，所以 \(k<\sqrt{k(k+1)}<k+\frac{1}{2}\). 对 \(k=1\)，\(2\)，\(\ldots\)，\(n\) 求和，得到 \(\frac{n(n+1)}{2}<a_n<\frac{n(n+2)}{2}<\frac{(n+1)^2}{2}\)，故所证不等式对所有正整数 \(n\) 成立.
\item 由第一问及 \(b_n=\frac{a_n}{n(n+1)}\)，有 \(\frac{1}{2}<b_n<\frac{n+1}{2n}=\frac{1}{2}+\frac{1}{2n}\). 任给 \(\varepsilon>0\)，取正整数 \(N=\left\lfloor\frac{1}{2\varepsilon}\right\rfloor+1\). 当 \(n\geqslant N\) 时，\(0<\left|b_n-\frac{1}{2}\right|<\frac{1}{2n}\leqslant\frac{1}{2N}<\varepsilon\). 这正是数列极限的定义，所以 \(\displaystyle\lim_{n\to\infty}b_n=\frac{1}{2}\).
\end{enumerate}
\end{solution}
\end{problem}

\begin{problem}

设 \(a\)，\(b\) 是两个实数，

\(A = \left\{\left(x, y\right)\mid x = n, y = na + b, n \in \Z\right\}\)，

\(B = \left\{\left(x, y\right)\mid x = m, y = 3m^2 + 15, m \in \Z\right\}\)，

\(C = \left\{(x, y)\mid x^2 + y^2 \leqslant 144\right\}\) 是平面 \(xOy\) 内的点集合，讨论是否存在 \(a\) 和 \(b\) 使得

\begin{enumerate}
\item \(A \cap B \neq \varnothing\)（\(\varnothing\) 表示空集）；
\item \(\left(a,b\right)\in C\) 同时成立.
\end{enumerate}

\begin{solution}
若 \(A\cap B\neq\varnothing\)，则存在整数 \(n\) 和 \(m\)，使得 \((n,na+b)=(m,3m^2+15)\). 比较横坐标得 \(n=m\)，记这个整数为 \(k\)，于是
\[
ka+b=3k^2+15
\]
恒等式 \((k^2+1)(a^2+b^2)=(ka+b)^2+(a-kb)^2\) 给出
\[
a^2+b^2\geqslant\frac{(ka+b)^2}{k^2+1}=\frac{(3k^2+15)^2}{k^2+1}
\]
而 \((3k^2+15)^2-144(k^2+1)=9(k^2-3)^2>0\)，其中严格不等号成立是因为 \(k\) 为整数，不可能有 \(k^2=3\). 所以 \(a^2+b^2>144\)，这与 \((a,b)\in C\)，即 \(a^2+b^2\leqslant144\)，矛盾. 因此不存在同时满足这两个条件的实数 \(a\) 和 \(b\).
\end{solution}
\end{problem}

\section{附加题}

\begin{problem}

已知曲线 \(y=x^3-6x^2+11x-6\). 在它对应于 \(x\in\left[0,2\right]\) 的弧段上求一点 \(P\)，使得曲线在该点的切线在 \(y\) 轴上的截距为最小，并求出这个最小值.

\begin{solution}
设切点的横坐标为 \(x\in[0,2]\)，则 \(f(x)=x^3-6x^2+11x-6\)，且 \(f'(x)=3x^2-12x+11\). 切线方程为 \(y-f(x)=f'(x)(X-x)\)，令 \(X=0\)，得到该切线在 \(y\) 轴上的截距
\[
b(x)=f(x)-xf'(x)=-2x^3+6x^2-6
\]
其导数为 \(b'(x)=-6x^2+12x=6x(2-x)\geqslant0\)（\(0\leqslant x\leqslant2\)），所以 \(b(x)\) 在 \([0,2]\) 上递增，最小值在 \(x=0\) 处取得. 此时 \(P=(0,f(0))=(0,-6)\)，切线斜率为 \(f'(0)=11\)，切线为 \(y=11x-6\)，在 \(y\) 轴上的截距为 \(-6\).
\end{solution}
\end{problem}
